problem:  
Let \((M^n,g)\) be a closed Riemannian manifold with \(\mathrm{Ric}_g \ge (n-1)g\) and denote by \(\lambda_1\) the first nonzero eigenvalue of the Laplacian acting on smooth functions.  
Assume that \(\lambda_1\) is close to \(n\), and the corresponding eigenspace has dimension \(k \ge n+1\) (i.e., higher than the multiplicity in the sphere case).  

Is it possible that such a manifold \(M\) is not Gromov–Hausdorff close to the unit sphere \(S^n\)?  
Equivalently, does there exist (for fixed \(n\)) a sequence of \(n\)-manifolds with \(\mathrm{Ric} \ge (n-1)g_i\) and eigenvalue–multiplicity pair  
\[
(\lambda_1(g_i), \dim E_{\lambda_1(g_i)}) \to (n, n+k)
\]
that fails to converge (in the pointed Gromov–Hausdorff or measured Gromov–Hausdorff sense) to a standard spherical space form?  

In other words, does *multiplicity pinching* of the first eigenvalue together with almost-equality of \(\lambda_1\) force the manifold to be quantitatively close (in geometry and topology) to a spherical quotient, suggesting a higher-order Obata-type rigidity?

Why is it a "good" problem:  
This formulation targets an essentially *unexplored direction* beyond known Lichnerowicz–Obata stability results. While much is known about the stability of the *value* of the first eigenvalue (Colding, Petersen, Aubry, etc.), almost nothing is known about the stability of its **multiplicity structure**, which in the spherical case encodes deep symmetry (via the \((n+1)\)-dimensional eigenspace generated by eigenfunctions corresponding to linear functions on \(\mathbb{R}^{n+1}\)).  

The problem thus links spectral geometry (eigenvalue multiplicity), comparison geometry (Ricci curvature bounds), and metric limit theory (Gromov–Hausdorff compactness). It poses a quantitative rigidity question at the interface of analysis and geometry that is simple to state yet profoundly difficult to resolve. Establishing or refuting such a “multiplicity pinching implies near-spherical symmetry” principle could lead to new insights into the geometric and analytic structure of nearly symmetric spaces—a potential next-generation Obata-type rigidity theorem.